\documentclass[11pt]{article}
\usepackage{amsmath,amssymb,amsfonts}
\usepackage[left=2.5cm,right=2.5cm,top=2cm,bottom=2cm]{geometry}
\usepackage{booktabs}
\usepackage{hyperref}

\DeclareMathOperator{\Var}{Var}
\DeclareMathOperator{\Cov}{Cov}
\DeclareMathOperator{\Corr}{Corr}
\DeclareMathOperator{\SD}{SD}
\newcommand{\bX}{\mathbf{X}}
\newcommand{\bZ}{\mathbf{Z}}
\newcommand{\bR}{\mathbf{R}}
\newcommand{\bI}{\mathbf{I}}
\newcommand{\bone}{\mathbf{1}}
\newcommand{\bt}{\mathbf{t}}
\newcommand{\bU}{\mathbf{U}}
\newcommand{\bu}{\mathbf{u}}

\title{Notebook Foundations:\\
  Development of the Mixed Effects Framework\\
  for Run-In Power Analysis}
\author{R.G.\ Thomas}
\date{Transcribed from handwritten notes, Aug--Oct 2019}

\begin{document}
\maketitle
\tableofcontents


\section{General linear mixed model}
\label{sec:lmm}

\noindent
\textit{(From notes dated 8/31/19 and 9/6/19.)}

\medskip
\noindent
Goal: reproduce Equations 11 and 13 from Frost et al.\ (2008).

\subsection{Outcome model}

For participant $i = 1, \ldots, N$ at visit $j = 0, \ldots, K$
in group $g_i \in \{0, 1\}$:
\begin{equation}
  Y_{ij} = \alpha + \alpha_{0i} + (\beta_1 + \beta_{1i}
    + \beta_x g_i)\, t_j + e_{ij}
\end{equation}
where $\alpha_{0i} \sim N(0, \sigma_\alpha^2)$ is a random
intercept, $\beta_{1i} \sim N(0, \sigma_\beta^2)$ is a random
slope, and $e_{ij} \sim N(0, \sigma^2)$ independently.

In matrix form for participant $i$:
\begin{equation}
  Y_i = X_i \beta + Z_i \gamma_i + \epsilon_i
  \label{eq:lmm}
\end{equation}
with
\begin{equation}
  X_i = \begin{pmatrix}
    1 & t_1 & t_1 g_i \\
    1 & t_2 & t_2 g_i \\
    \vdots & \vdots & \vdots \\
    1 & t_K & t_K g_i
  \end{pmatrix}, \quad
  Z_i = \begin{pmatrix}
    1 & t_1 \\
    1 & t_2 \\
    \vdots & \vdots \\
    1 & t_K
  \end{pmatrix}, \quad
  \beta = \begin{pmatrix}
    \alpha \\ \beta_1 \\ \beta_x
  \end{pmatrix}, \quad
  \gamma_i = \begin{pmatrix}
    \alpha_{0i} \\ \beta_{1i}
  \end{pmatrix}
\end{equation}


\subsection{Marginal distribution}

Since $\gamma_i \sim N(\mathbf{0}, D)$ and
$\epsilon_i \sim N(\mathbf{0}, R_i)$ independently:
\begin{align}
  E[Y_i \mid \gamma_i] &= X_i\beta + Z_i\gamma_i, \qquad
    V(Y_i \mid \gamma_i) = R_i \\
  Y_i &\sim N(X_i\beta,\; \Sigma_i)
\end{align}
where the marginal covariance is
\begin{equation}
  \Sigma_i = V(Y_i)
  = V(X_i\beta) + V(Z_i\gamma_i) + V(\epsilon_i)
  = 0 + Z_i D Z_i' + R_i
  = R_i + Z_i D Z_i'
  \label{eq:sigma}
\end{equation}


\subsection{Simplified model (random slope only)}

Dropping the random intercept
($\sigma_\alpha^2 = 0$, $D = \sigma_\beta^2$ scalar):
\begin{equation}
  Y_{ij} = \beta + (\beta_x I_{x,i} + u_{1,i})\, t_j
    + \epsilon_{ij}
\end{equation}
with $u_{1,i} \sim N(0, \sigma_1^2)$,
$\epsilon_{ij} \sim N(0, \sigma^2)$. Then:
\begin{equation}
  X_i = \begin{pmatrix}
    t_0 & t_0 g_i \\
    \vdots & \vdots \\
    t_K & t_K g_i
  \end{pmatrix}, \quad
  U_i = \begin{pmatrix}
    t_0 \\ \vdots \\ t_K
  \end{pmatrix}, \quad
  G = (\sigma_b^2), \quad
  R = \sigma^2 \bI_{K+1}
\end{equation}
and $V_i = R + U_i G U_i' = \sigma^2 \bI + \sigma_b^2
\bu\bu'$.


\section{Variance of the treatment effect with intercepts}
\label{sec:var_full}

\noindent
\textit{(From notes dated 9/6/19, page 13.)}

\medskip
\noindent
With both random intercept and random slope, the variance
of the treatment effect estimator is
\begin{equation}
  \Var(\hat\beta_{TX})
  = \frac{2}{m}\left[
    \frac{\sigma^4 / K + \sigma^2(\sigma_\beta^2
      + \sigma_\alpha^2)
      + \sigma_\alpha^2 \sigma_\beta^2(1 - \rho^2)}
    {K[\sigma^2 + K\sigma_\alpha^2]}
  \right]
  \label{eq:var_full}
\end{equation}
where $\rho = \sigma_{12}/(\sigma_1 \sigma_2)$ is the
correlation between random intercepts and slopes.

This simplifies to
\begin{equation}
  \frac{2}{m}\left[
    \frac{\sigma^2/K + \sigma_\beta^2 + \sigma_\alpha^2
      + \frac{\sigma_\alpha^2 \sigma_\beta^2}{\sigma^2}
        (1-\rho^2)}
    {1 + K\frac{\sigma_\alpha^2}{\sigma^2}}
  \right]
\end{equation}

When $\sigma_\alpha^2 = 0$ (no random intercept), this
reduces to
\begin{equation}
  \Var(\hat\beta_{TX})
  = \frac{2}{m}\left[
    \frac{\sigma^2}{K} + \sigma_\beta^2
  \right]
\end{equation}
which is the Frost result for the conventional design
written in terms of the number of post-baseline visits $K$.


\section{Change-score parameterization}
\label{sec:changescore}

\noindent
\textit{(From notes dated 9/25/19, page 17.)}

\medskip
\noindent
The full model is
\begin{equation}
  y_{ij} = \alpha_0' + \alpha_{0i}'
    + (\alpha_1' + \alpha_{1i}' + \alpha_x' g_i)\,t_j
    + e_{ij}
\end{equation}

Change scores relative to baseline at $j = 0$:
\begin{align}
  z_{ij} &= y_{ij} - y_{i0} \notag \\
  &= (\alpha_1' + \alpha_{1i}' + \alpha_x' g_i)(t_j - t_0)
    + (e_{ij} - e_{i0}) \notag \\
  &= (\alpha_1' + \alpha_{1i}' + \alpha_x' g_i)\,s_j
    + \varepsilon_j
  \label{eq:change}
\end{align}
where $s_j = t_j - t_0$ and
$\varepsilon_j \sim N(0, \sigma^2 \bI)$ for
$j = 1, \ldots, K$.

In matrix form:
\begin{equation}
  z_i = X\beta + U b_i + \varepsilon_i
\end{equation}
with $b_i \sim N(0, D)$,
$\varepsilon_i \sim N(0, R)$.

The change-score parameterization eliminates the intercept
$\alpha_0'$ and the random intercept $\alpha_{0i}'$.


\section{Normalized time and the algorithmic recipe}
\label{sec:algorithm}

\noindent
\textit{(From notes dated 9/25/19, page 16.)}

\subsection{Time normalization}

Define normalized time:
\begin{equation}
  t_i^* = \frac{t_i - \bar{t}}{\SD(t_i)}
\end{equation}
so that $\sum t_i^* = 0$, $\sum t_i^{*2} = K - 1$,
and $V(t_i^*) = 1$. The original-scale variance is
recovered via
\begin{equation}
  V(\hat\alpha_x) = V(\hat\alpha_x^*)
    \cdot \frac{1}{s_x^2}
\end{equation}
where $s_x^2 = \Var(t_i)$.

\subsection{Seven-step algorithm}

For a two-arm trial with $m$ treatment and $m$ placebo
participants:

\begin{enumerate}
  \item Compute the per-subject marginal covariance:
    $V_i^* = R_i + Z_i^* D\, Z_i^{*\prime}$

  \item Invert: $V_i^{*-1}$

  \item Compute $X_i' V_i^{-1} X_i$ for treatment
    ($g_i = 1$)

  \item Compute $X_i' V_i^{-1} X_i$ for placebo
    ($g_i = 0$)

  \item Sum the information matrix:
    $F = \sum_{\text{trt}} X_i' V_i^{-1} X_i
      + \sum_{\text{plc}} X_i' V_i^{-1} X_i$

  \item Invert: $F^{-1}$

  \item Extract: $F_{2,2}^{-1} = V(\hat\alpha_x^*)$;
    then $V(\hat\alpha_x) = V(\hat\alpha_x^*) / s_x^2$
\end{enumerate}

The Woodbury identity is applied at Step~2 to avoid
direct inversion of $V_i^*$.


\section{Woodbury identity for the raw-outcome model}
\label{sec:woodbury_raw}

\noindent
\textit{(From notes dated 9/25/19, pages 14--15.)}

\subsection{Setup}

With $R = \sigma^2 \bI_{K+1}$ and scalar $G = \sigma_b^2$:
\begin{equation}
  V = UGU' + R = \sigma_b^2 \bu\bu' + \sigma^2 \bI
\end{equation}

Applying Woodbury with $A = \sigma^2 \bI$,
$C = \sigma_b^2$:
\begin{align}
  V^{-1} &= \frac{1}{\sigma^2}\bI
    - \frac{1}{\sigma^4}\,\bu
    \left(\frac{1}{\sigma_b^2}
      + \frac{1}{\sigma^2}\bu'\bu\right)^{-1}
    \bu' \notag \\
  &= \frac{1}{\sigma^2}\left[
    \bI - \frac{\sigma_b^2}{\sigma^2}\,\bu\bu'
  \right]
  \label{eq:Vinv_raw}
\end{align}
where the last line holds when $\sum t_i = 0$ (which
makes $\bu'\bu = \sum t_i^2$ and the cross terms vanish).


\subsection{Information matrix}

The total information is
\begin{equation}
  X'\Sigma^{-1}X = \sum_{\text{trt}} X_i' V^{-1} X_i
    + \sum_{\text{plc}} X_i' V^{-1} X_i
  = 2m\, X_i' V^{-1} X_i^{(g=1)}
    + m\, X_i' V^{-1} X_i^{(g=0)}
\end{equation}

With $t_K = Kd$ (equally spaced), the outer product
$\bu\bu'$ has elements
\begin{equation}
  [\bu\bu']_{jk} = t_j t_k = d^2 jk
  \qquad\Longrightarrow\qquad
  \bu\bu' = d^2 \begin{pmatrix}
    1 & 2 & \cdots & K \\
    2 & 4 & \cdots & 2K \\
    \vdots & & \ddots & \vdots \\
    K & 2K & \cdots & K^2
  \end{pmatrix}
\end{equation}

The Woodbury correction to the information matrix is
\begin{equation}
  X'R^{-1}Z \cdot (G^{-1} + Z'R^{-1}Z)^{-1}
    \cdot Z'R^{-1}X
\end{equation}
which is computed per-group and then summed, because
$X$ differs between treatment and placebo.


\section{Frost 3.1.1 in the change-score parameterization}
\label{sec:frost311}

\noindent
\textit{(From ``Take 2'' notes, pages 9--10.)}

\subsection{Design matrices}

For the conventional design ($r = 0$, $k = 2$) with change
scores:
\begin{align}
  X^{T} &= t\begin{pmatrix}
    1 & 0 & 0 \\ 0 & 1 & 1
  \end{pmatrix}, &
  X^{P} &= t\begin{pmatrix}
    1 & 0 & 0 \\ 0 & 1 & 0
  \end{pmatrix} \\[4pt]
  Z &= t\begin{pmatrix} 1 \\ 1 \end{pmatrix}, &
  G &= b, \quad
  R = a\begin{pmatrix} 2 & -1 \\ -1 & 2 \end{pmatrix}
\end{align}

Note: the ``Take 2'' version uses a 3-parameter design
($\delta, \beta, \gamma$) even for the conventional case,
treating the first observation as a separate phase. This
is equivalent to Frost 3.1.2 applied to the $r=0$ case.

\subsection{Intermediate computations}

\begin{equation}
  R^{-1} = \frac{1}{3a}\begin{pmatrix}
    2 & 1 \\ 1 & 2
  \end{pmatrix}
\end{equation}

\textbf{$F = (G^{-1} + Z'R^{-1}Z)^{-1}$:}
\begin{align}
  Z'R^{-1}Z &= t^2 (1,\,1) \cdot \frac{1}{3a}
    \begin{pmatrix} 2&1\\1&2 \end{pmatrix}
    \begin{pmatrix} 1\\1 \end{pmatrix}
  = \frac{t^2}{3a}(1,\,1)\begin{pmatrix}3\\3\end{pmatrix}
  = \frac{6t^2}{3a}
  = \frac{2t^2}{a}
\end{align}
\begin{equation}
  \boxed{
    F = \left(\frac{1}{b}
      + \frac{2t^2}{a}\right)^{-1}
    = \frac{ab}{a + 2bt^2}
  }
\end{equation}

\textbf{$X'R^{-1}X$ for treatment:}
\begin{equation}
  X^{T\prime} R^{-1} X^T
  = \frac{t^2}{3a}\begin{pmatrix}
    1&0\\0&1\\0&1
  \end{pmatrix}
  \begin{pmatrix}2&1\\1&2\end{pmatrix}
  \begin{pmatrix}1&0&0\\0&1&1\end{pmatrix}
  = \frac{t^2}{3a}\begin{pmatrix}
    2 & 1 & 1 \\ 1 & 2 & 2 \\ 1 & 2 & 2
  \end{pmatrix}
\end{equation}

\textbf{$X'R^{-1}X$ for placebo:}
\begin{equation}
  X^{P\prime} R^{-1} X^P
  = \frac{t^2}{3a}\begin{pmatrix}
    2 & 1 & 0 \\ 1 & 2 & 0 \\ 0 & 0 & 0
  \end{pmatrix}
\end{equation}

\textbf{Sum:}
\begin{equation}
  X'R^{-1}X
  = \frac{t^2}{3a}\begin{pmatrix}
    4 & 2 & 1 \\ 2 & 4 & 2 \\ 1 & 2 & 2
  \end{pmatrix}
\end{equation}

\textbf{$X'R^{-1}Z$ for treatment:}
\begin{equation}
  X^{T\prime} R^{-1} Z
  = \frac{t^2}{3a}\begin{pmatrix}1&0\\0&1\\0&1\end{pmatrix}
    \begin{pmatrix}2&1\\1&2\end{pmatrix}
    \begin{pmatrix}1\\1\end{pmatrix}
  = \frac{t^2}{3a}\begin{pmatrix}3\\3\\3\end{pmatrix}
  = \frac{t^2}{a}\begin{pmatrix}1\\1\\1\end{pmatrix}
\end{equation}

\textbf{$X'R^{-1}Z$ for placebo:}
\begin{equation}
  X^{P\prime} R^{-1} Z
  = \frac{t^2}{a}\begin{pmatrix}1\\1\\0\end{pmatrix}
\end{equation}

\textbf{$Z'R^{-1}Z$:}
\begin{equation}
  Z'R^{-1}Z = \frac{2t^2}{a}
\end{equation}

\subsection{Woodbury correction}

\begin{align}
  W &= X^{T\prime}R^{-1}Z \cdot F \cdot Z'R^{-1}X^T
    + X^{P\prime}R^{-1}Z \cdot F \cdot Z'R^{-1}X^P
    \notag \\
  &= \frac{t^4}{a^2}\cdot\frac{ab}{a+2bt^2}\left[
    \begin{pmatrix}1\\1\\1\end{pmatrix}(1,1,1)
    + \begin{pmatrix}1\\1\\0\end{pmatrix}(1,1,0)
  \right] \notag \\
  &= \frac{t^4 b}{a(a+2bt^2)}\begin{pmatrix}
    2 & 2 & 1 \\ 2 & 2 & 1 \\ 1 & 1 & 1
  \end{pmatrix}
\end{align}

\subsection{Information matrix and variance}

\begin{equation}
  X'\Sigma^{-1}X = X'R^{-1}X - W
  = \frac{t^2}{3a}\begin{pmatrix}
    4&2&1\\2&4&2\\1&2&2
  \end{pmatrix}
  - \frac{t^4 b}{a(a+2bt^2)}\begin{pmatrix}
    2&2&1\\2&2&1\\1&1&1
  \end{pmatrix}
\end{equation}

Factoring:
\begin{equation}
  X'\Sigma^{-1}X = \frac{t^2}{a}\left[
    \frac{1}{3}\begin{pmatrix}
      4&2&1\\2&4&2\\1&2&2
    \end{pmatrix}
    - \frac{t^2 b}{a+2bt^2}\begin{pmatrix}
      2&2&1\\2&2&1\\1&1&1
    \end{pmatrix}
  \right]
  \label{eq:XSigX_frost}
\end{equation}

Define $d = \frac{t^2 b}{a + 2bt^2}$. The $(3,3)$ element
of $(X'\Sigma^{-1}X)^{-1}$ gives $\Var(\hat\gamma)$.


\section{AR(1) covariance sums}
\label{sec:ar1}

\noindent
\textit{(From notes dated 10/14, pages 2--7.)}

\medskip
\noindent
For an AR(1) correlation structure with parameter $c = \rho$,
the variance of the mean change score involves sums of the
form
\begin{equation}
  S_n = \sum_{i=1}^{n-1} (n - i)\, c^{i-1}
\end{equation}

\subsection{Closed-form derivation}

Rewriting:
\begin{align}
  S_n &= \sum_{i=0}^{n-1} (n-i)\, c^i - n \notag \\
  &= c^0 \sum_{i=0}^{n-1} n\, c^i
    - c^0 \sum_{i=0}^{n-1} i\, c^i - n
\end{align}

Using the geometric series
$\sum_{i=0}^{n-1} c^i = \frac{1-c^n}{1-c}$
and the derivative identity
$\sum_{i=0}^{n-1} i\, c^i
= \frac{c - nc^n + (n-1)c^{n+1}}{(1-c)^2}$:

\begin{align}
  S_n &= n\cdot\frac{1-c^n}{1-c}
    - \frac{c - nc^n + (n-1)c^{n+1}}{(1-c)^2} - n
    \notag \\
  &= \frac{n(1-c^n)(1-c) - c + nc^n - (n-1)c^{n+1}
    - n(1-c)^2}{(1-c)^2}
\end{align}

Expanding and collecting:
\begin{equation}
  \boxed{
    S_n = \frac{n - 2nc + nc^2 - c + c^n(n-1)(c-1)
      + c^{n+1}}{(1-c)^2}
  }
\end{equation}

\subsection{Verification for small $n$}

\textbf{$n = 2$:} $S_2 = 1$.
\begin{equation}
  \frac{2 - 4c + 2c^2 - c + c^2(c - 1) + c^3}{(1-c)^2}
  = \frac{2 - 5c + 3c^2 + c^3 - c^2}{(1-c)^2}
  = \frac{(1-c)^2}{(1-c)^2} = 1 \quad\checkmark
\end{equation}

\textbf{$n = 3$:} $S_3 = 2 + c$.
\begin{equation}
  \frac{3 - 6c + 3c^2 - c + 2c^3(c-1) + c^4}{(1-c)^2}
  = \frac{3 - 6c + 6c^2 - 3c^3 + c^4 - c}{(1-c)^2}
\end{equation}
which simplifies (via pages 5--6) to $2 + c$. $\checkmark$

\textbf{$n = 4$:} $S_4 = 3 + 2c + c^2$.
Verified on page 4 via the identity
$\sum_{i=1}^{3}(4-i)c^{i-1} = 3 + 2c + c^2$.
$\checkmark$

\textbf{$n = 5$:} $S_5 = 4 + 3c + 2c^2 + c^3$.
Verified by factoring: $(1-c^5)/(1-c) = 1 + c + c^2 + c^3
+ c^4$ and checking the polynomial identity. $\checkmark$

\subsection{Alternative form}

For the special case where $S_n$ appears as a coefficient
in the variance formula, the sum admits the factored form
(page 5):
\begin{equation}
  S_n = \frac{1}{1-c}\left[
    n + \frac{1 - c^n}{1 - c}
  \right]
  = \frac{1}{1-c}\left[
    n + 1 + c + c^2 + \cdots + c^{n-1}
  \right]
\end{equation}

This can be rewritten as (page 7):
\begin{equation}
  S_n = \frac{(c^{n-1})(1-c) + (c^n - 1)}{(1-c)^2}
  = \frac{c^n - 1 + n - nc}{(c - 1)^2}
\end{equation}

The final simplified form (page 7, sideways) is
\begin{equation}
  \boxed{
    S_n = 1 + c + c^2 + \cdots + c^{n-1} - n
    \cdot\frac{c - 1}{c - 1}
    = \frac{n(1-c) - (1 - c^n)}{(1-c)^2}
  }
\end{equation}


\section{Stratified variance under unequal visit patterns}
\label{sec:stratified}

\noindent
\textit{(From notes dated 10/14, page 1.)}

\medskip
\noindent
When individuals have different numbers of run-in visits
$r_i$ and final visits $f_i$, the variance is computed by
stratification. Let $z_{ij} = X_{ijk} - X_{ij0}$ denote
the change score. The treatment effect variance is
\begin{equation}
  V(\bar{z}_1 - \bar{z}_2)
  = V(\bar{z}_1) + V(\bar{z}_2)
  = \frac{1}{m_1^2}\,V\!\left(
    \sum_{i=1}^{m_1}(X_{i1k} - X_{i10})\right)
  + \frac{1}{m_2^2}\,V\!\left(
    \sum_{i=1}^{m_2}(Y_{i2k} - Y_{i20})\right)
\end{equation}

Stratify participants by visit pattern
$(r_i, f_i)$:
\begin{equation}
  \sum_{i=1}^{m} V(X_{ik} - X_{i0})
  = \sum_{s=1}^{S}\; \sum_{i \in \text{stratum } s}
    V(X_{ik} - X_{i0})
\end{equation}
where $S$ is the number of distinct strata.

Distribution of strata: $p_1 \% \times p_2 \%$.

Within each stratum (with $r$ run-in, $f$ final visits):
\begin{align}
  V(X_{ik} - X_{i0})
  &= V\!\left(
    \tfrac{1}{c}\sum F_i - \tfrac{1}{r}\sum R_i
  \right) \notag \\
  &= V\!\left(\tfrac{1}{c}\sum F_i\right)
    + V\!\left(\tfrac{1}{r}\sum R_i\right)
    - 2\,\Cov\!\left(
      \tfrac{1}{c}\sum F_i,\;
      \tfrac{1}{r}\sum R_i
    \right)
\end{align}
where $F_i$ and $R_i$ denote the final (post-randomization)
and run-in observations, respectively.


\section{Appendix: Frost 3.1.1 via direct inversion}
\label{sec:direct}

\noindent
\textit{(From notes dated 8/25/19, ``Take 2'' pages 9--10.)}

\medskip
\noindent
As a cross-check, the Frost 3.1.1 result can be obtained
by directly computing $(G^{-1} + Z'R^{-1}Z)^{-1}$ without
the Woodbury identity.

With $Z = t(1,\, 1)'$ and
$R^{-1} = \frac{1}{3a}\begin{pmatrix}2&1\\1&2\end{pmatrix}$:

\begin{equation}
  G^{-1} + Z'R^{-1}Z
  = \frac{1}{b} + \frac{t^2}{3a}(1,\,1)
    \begin{pmatrix}2&1\\1&2\end{pmatrix}
    \begin{pmatrix}1\\1\end{pmatrix}
  = \frac{1}{b} + \frac{6t^2}{3a}
  = \frac{a + 2bt^2}{ab}
\end{equation}
so
\begin{equation}
  H = (G^{-1} + Z'R^{-1}Z)^{-1} = \frac{ab}{a + 2bt^2}
\end{equation}

This is the scalar $H$ that appears in the Woodbury
correction, confirming the ``Take 2'' derivation.

The final variance (from the report, verified here):
\begin{equation}
  \Var(\hat\gamma)\big|_{r=0,\,k=2}
  = \frac{a + 2bt^2}{t^2}
  = \frac{\sigma^2 + 2\sigma_b^2 t^2}{t^2}
\end{equation}

\end{document}
